\documentclass[a4paper,conference]{IEEEtran}

\usepackage{graphicx}
\usepackage{booktabs}
\usepackage{tabularx}
\usepackage{cite}
\usepackage{url}
\usepackage[hidelinks]{hyperref}
\usepackage{dblfloatfix}    % fix ordering so figure* never lands after the references

% ---- float behaviour: keep wide figures near their text, never dumped at the end ----
\setcounter{topnumber}{3}
\setcounter{dbltopnumber}{3}
\setcounter{totalnumber}{5}
\renewcommand{\topfraction}{0.95}
\renewcommand{\dbltopfraction}{0.95}
\renewcommand{\bottomfraction}{0.6}
\renewcommand{\textfraction}{0.05}
\renewcommand{\floatpagefraction}{0.75}
\renewcommand{\dblfloatpagefraction}{0.75}
\hyphenation{multi-lin-gual multi-func-tion-al-ity multi-gran-u-lar-ity re-triev-er re-triev-ers}

% ---------------------------------------------------------------------------
% AUTHOR: fill these six placeholders, then delete this block.
% Five come from analysis/transfer_stats.py (stats.json); one is the comparator.
% Until filled they render as visible [TODO: ...] markers in the PDF.
% ---------------------------------------------------------------------------
\newcommand{\TODO}[1]{\textbf{[TODO: #1]}}
\newcommand{\wtPatCI}{[-0.005, 0.006]}
\newcommand{\wtArcCI}{[-0.005, 0.011]}
\newcommand{\wtQwenCI}{[-0.006, 0.009]}
\newcommand{\maxArgmaxProb}{0.68}
\newcommand{\ndcgtenCI}{[0.053, 0.074]}
\newcommand{\selectedSpec}{ARM-03: datalyes/patembed-large (revision 2d5c0f92a3e5dc3d5415c08e612c57543c0e03ad); a2-arm-03-matched-b2-orthogonal, title+abstract+claims, 384-token passages with 64-token overlap, NFKC whitespace, mean non-padding pooling, L2 normalization, max-p family aggregation; query prefix ``encode query for different document retrieval:''; no fusion, top-100.}
\newcommand{\comparatorSpec}{FAST: BM25S 0.3.10 plus Snowflake/snowflake-arctic-embed-m-v2.0 (revision 95c2741480856aa9666782eb4afe11959938017f); P00-TAC-DOC title+abstract+claims as one family document with NFKC canonical-whitespace, case-preserving construction; dense query prefix ``query:''; synchronous reciprocal-rank fusion (RRF-60), top-100.}

\title{Beyond the Retriever: Representation Transfer in Cross-Domain Patent Retrieval}

% Double-anonymous review version. CAMERA-READY: replace with real author block.
\author{Anonymous Authors}

\begin{document}
\maketitle

\begin{abstract}
A patent retriever never sees the invention it is asked to find. It sees a representation we build for it---a title, a claim, a window of passages---and scores that. Cross-domain benchmarks freeze this representation and read the resulting differences as facts about the retrievers, as if the construction were a neutral control. It is not: it is part of the system under test. We take that assumption apart on family-level cross-domain patent retrieval (DAPFAM) by freezing five retrievers and varying five deterministic constructions, keeping development, one-time selection, and held-out confirmation as separate evidence levels. The result is deflationary. Across a $3\times3$ cross-retriever transfer matrix, the best construction changes from one target retriever to the next---but the changes sit inside bootstrap noise (within-target Recall@100 spread below 0.004), while the gap between retrievers is an order of magnitude larger. At the coarse, field level, representation does not reorder retrievers; retriever identity does. A preregistered complete configuration then wins cleanly on 872 held-out queries, lifting Recall@100 from 0.331 to 0.442 (paired difference 0.111, 95\% bootstrap CI [0.102, 0.120]) over a frozen comparator---a system-level result our design does not credit to representation alone. Applied unchanged to the full benchmark, that winner reveals what actually caps recall: 78\% of relevant families never enter the Top-200 pool, and a perfect reordering of what remains would add only 0.072. Cross-domain patent retrieval is exposure-bound, not ordering-bound. We argue that retrieval quality should be reported at the retriever--representation configuration level, and that effort spent tuning the ranker is wasted while most of the evidence is missing from the pool.
\end{abstract}

\begin{IEEEkeywords}
patent retrieval, cross-domain retrieval, document representation, dense retrieval, information retrieval
\end{IEEEkeywords}

\section{Introduction}

A patent retriever never sees a patent. It sees the representation we choose to hand it---a title and abstract, a full claim set, a single structured claim, a window of passages. Index a lexical and a dense retriever on the same construction and they appear to receive the same evidence, so benchmarks freeze the construction and compare the models. The move is convenient, and it hides an assumption: that a fixed representation is a control the model cannot see around, rather than one more part of the system being measured. Freezing a variable does not neutralize it. It only hides it. So we ask a blunt question: with the retriever fixed, does the representation still move the ranking?

Patent families make the question concrete. The same invention can be exposed as a bare title, a single structured claim, overlapping passages, or a fused set of section views, and earlier work has shown that granularity and field choice change effectiveness, most sharply across technological domains~\cite{ayaou2026dapfam,yousefiramandi2026patent}. What that work leaves open is transfer: if a construction is tuned around one retriever, does its advantage survive when a different retriever consumes it?

\begin{figure*}[!t]
\centering
\includegraphics[width=0.98\textwidth]{figures/rebuilt/overview_evidence_map.pdf}
\caption{\textbf{Evidence map.} Representation construction and transfer analysis are development evidence. Development closes before Selection-125 and the Final-872 complete-system confirmation; the full-benchmark diagnosis is post-confirmatory. Each stage answers a different question, and none borrows another's authority.}
\label{fig:evidence-map}
\end{figure*}

We freeze five heterogeneous retrievers and vary five deterministic constructions, with no fine-tuning, adapters, or distillation, then run a full $3\times3$ transfer matrix across three dense retrievers and ask, plainly, \emph{best representation---for whom?} We expected the answer to reorder the field. It does not. The nominally best construction flips from target to target, but the flips are smaller than the noise around them, while the retrievers stay separated by an order of magnitude. At the coarse level of field selection, representation is not the lever; retriever identity is. This is a claim about coarse constructions, not about representation in general: richer, learned representation programs can move even dense retrievers substantially~\cite{onuallain2026autoindex}, a boundary we draw carefully in Section~\ref{sec:related}.

What makes these claims worth trusting is discipline, so we keep three kinds of evidence apart and never let one speak with another's authority (Figure~\ref{fig:evidence-map}). Development characterizes representation response and transfer. A one-time selection stage then freezes a complete configuration, which a protected 872-query held-out comparison confirms. Only after that system is fixed do we turn its full-benchmark Top-200 pool into a diagnostic---and it delivers the paper's sharpest finding: recall is capped not by how we order candidates, but by which candidates ever appear. Most relevant families never reach the pool. Our contributions are (i)~controlled multi-retriever evidence that coarse representation choices do not transfer into retriever reordering, (ii)~a preregistered, held-out confirmation of a complete configuration, and (iii)~a post-confirmatory diagnosis showing that cross-domain patent retrieval is exposure-bound.

\section{Related Work}
\label{sec:related}

Patent retrieval benchmarks increasingly expose the effects of domain and document construction. DAPFAM gives family-level in-domain and out-of-domain evaluation and systematically compares lexical and dense retrieval, document and passage units, aggregation, and fusion~\cite{ayaou2026dapfam}. PatenTEB broadens patent embedding evaluation across retrieval and other patent-language tasks, and a recent 22-model benchmark examines multiple text-section views and reports stubborn out-of-domain degradation~\cite{ayaou2025patenteb,yousefiramandi2026patent}. Together they establish that model family and document view both matter. Much of this is one research line (the DAPFAM, PatenTEB, and PatEmbed resources share authorship), and our question is orthogonal to it: not whether a view helps a given retriever, but whether a view chosen around one retriever stays best when another retriever consumes it, under a protocol that keeps development, selection, and confirmation apart.

The sharpest contrast is with representation-optimization work, and because it cuts against our headline we meet it directly. AutoIndex learns executable representation programs---slicing, enriching, reweighting, normalizing documents---around a fixed BM25 backend, and reports large recall gains from representation alone, including a preliminary dense result in which a learned program improves Qwen3-Embedding-0.6B Recall@100 by 18.3\% on a single task~\cite{onuallain2026autoindex}. Our study is complementary and deliberately narrower. We hold five retrievers fixed and vary a small set of coarse, field-level static constructions; we do not search a rich program space. So our finding that these constructions do not reorder retrievers is a statement about the portability of coarse choices across frozen dense retrievers, not a claim that representation is inert---AutoIndex shows the opposite once the representation space is rich and corpus-adaptive. Seen this way, our transfer matrix is an early empirical probe of exactly the cross-retriever transfer question AutoIndex leaves open.

Finally, multi-stage retrieval separates candidate generation from downstream scoring: a cheap first stage supplies a bounded pool, and more expensive models rerank it~\cite{asadi2013candidate,nogueira2019passage}. We invoke that distinction only after confirmation, to measure candidate exposure and to compute a perfect-ordering bound inside the fixed Top-200 pool---separating error a reranker could still recover from error that is simply absent.

\section{Study Design}

\textbf{Benchmark and systems.} The evaluation unit is a patent family. The benchmark holds 1,247 query families, 45,336 target families, and 49,869 judged relations. Cross-domain Recall@100 is primary; nDCG@100 and nDCG@10 are secondary~\cite{jarvelin2002cumulated}. We freeze five retrieval systems: BM25~\cite{robertson2009probabilistic}, BGE-M3~\cite{chen2024bgem3}, PatEmbed-large~\cite{ayaou2025patenteb}, Arctic Embed M v2.0~\cite{yu2024arctic}, and Qwen3 Embedding 0.6B~\cite{zhang2025qwen3}. No system receives fine-tuning, adapters, distillation, or continued pretraining during the study.

\begin{table}[!t]
\caption{Separated evidence roles and query populations. The full-benchmark strict cross-domain population is not the Final-872 population.}
\label{tab:protocol}
\centering
\footnotesize
\setlength{\tabcolsep}{3.5pt}
\begin{tabular}{@{}lrl@{}}
\toprule
Evidence role & Queries & Purpose \\
\midrule
Representation development & 150 & Common/per-system search \\
Combined development & 250 & Transfer and composition \\
One-time selection & 125 & Freeze final systems \\
Protected confirmation & 872 & Complete-system comparison \\
Full benchmark & 1,247 & Characterization/diagnosis \\
\bottomrule
\end{tabular}
\end{table}

\textbf{Representations and evidence separation.} The shared deterministic surface holds five constructions: title--abstract--claims as one family document; title and abstract; the first structured independent claim; 384-token passages with 64-token overlap and family-level maximum aggregation; and labeled title, abstract, and claims views combined by family-level rank fusion~\cite{cormack2009rrf}. Table~\ref{tab:protocol} lists the protected sequence. Representation development uses 150 queries; transfer and fixed-composition controls cover all 250 development queries. Four profiles pass through one Selection-125 exposure before the research system and a static comparator are frozen. Final confirmation uses paired query-level differences, 10,000 paired bootstrap resamples, and 95\% percentile intervals~\cite{davison1997bootstrap}. After confirmation, the winner alone is applied to all 1,247 queries unchanged except retrieval depth, extended to 200 for the diagnostic pool; every later analysis reuses that one Top-200 pool with no further retrieval or model selection.

\section{Cross-Retriever Representation Transfer}

\begin{figure*}[!t]
\centering
\includegraphics[width=0.98\textwidth]{figures/rebuilt/fig1_a3_transfer.pdf}
\caption{\textbf{Best representation---for whom?} Development Recall@100 for three representation sources and consuming retrievers, on disclosed target-specific values. The nominal best source (amber) changes with the target, but the within-target spread is below 0.004 for every target, while the between-target bands differ by an order of magnitude. Development evidence; no single cell carries confirmatory inference.}
\label{fig:transfer}
\end{figure*}

\emph{If representation were neutral preprocessing, frozen retrievers would respond to it alike.} We test that first on a common screen: five systems under five shared constructions, 25 development cells. PatEmbed, Arctic Embed, and Qwen3 carry forward to the transfer analysis; BM25 and BGE-M3 stay diagnostic. The screen probes one predefined surface on one benchmark, and is not a model leaderboard.

A per-system search then registers 52 configurations in advance. Registering the whole set---including the eight conditional reserves whose activation predicates never fired, leaving 44 executed---fixes the search space before any score is seen, which is the point of registering it. The responses are mixed: Arctic improves strictly under its rule, PatEmbed produces a within-search winner that ties its incumbent, Qwen3 shows no strict improvement, and BM25 and BGE-M3 end in diagnostic ties. Because the common screen and the per-system search use different decision rules, their numbers are not a before/after pair; Table~\ref{tab:development-response} reports each with its frozen interpretation.

\begin{table}[!t]
\caption{Development response across frozen retrievers. Shared-screen and per-system-search Recall@100 are distinct summaries produced under different decision rules; their difference is not a causal before/after effect.}
\label{tab:development-response}
\centering
\footnotesize
\setlength{\tabcolsep}{2.1pt}
\begin{tabular}{@{}lccp{0.40\columnwidth}@{}}
\toprule
Retriever & \shortstack{Shared\\screen} & \shortstack{Per-system\\search} & Frozen interpretation \\
\midrule
BM25 & 0.191 & 0.235 & Diagnostic tie; no winner \\
BGE-M3 & 0.270 & 0.290 & Diagnostic tie; no winner \\
PatEmbed & 0.413 & 0.423 & Numerical tie to incumbent; retained \\
Arctic Embed & 0.341 & 0.359 & Strict improvement; retained \\
Qwen3 Embed. & 0.364 & 0.374 & No strict improvement; transfer control \\
\bottomrule
\end{tabular}
\end{table}

The mixed responses leave a tempting possibility open: a construction discovered around one system might still transfer to the others. We test it with a complete $3\times3$ matrix---each construction selected from PatEmbed, Arctic Embed, or Qwen3, consumed by each of the three targets, every cell over all 250 development queries.

Figure~\ref{fig:transfer} lays the result out. Read target by target, the nominally best construction does move: Qwen3-derived text wins for PatEmbed, Arctic-derived for Arctic, PatEmbed-derived for Qwen3, and the nDCG@100 ordering disagrees again (its top cell is Arctic-derived text on PatEmbed, 0.352). No construction holds first place everywhere. But the differences behind that reordering are tiny, and because every cell shares the same 250 queries we can test them head-on. Within each target, the paired gap between the best and second-best construction has a 95\% bootstrap interval that includes zero (\wtPatCI\ for PatEmbed, \wtArcCI\ for Arctic, \wtQwenCI\ for Qwen3), and no construction is a stable within-target winner under resampling (best-source probability at most \maxArgmaxProb). The reordering is noise. What is not noise is the retriever: every construction lands near 0.419 Recall@100 on PatEmbed, near 0.34 on Arctic, near 0.36 on Qwen3---bands separated by far more than the spread inside any one of them. Change the construction and you shuffle positions within a band; change the retriever and you move between bands.

A composition control tells the same story from another angle. Two-system fusion (0.419) barely edges the strongest single system (0.418), and three-system fusion drops to 0.415: stacking retrievers does not monotonically help.

At the coarse, field level, representation is a passenger, not the driver; no single transfer cell carries confirmatory weight. That settles development and sets up a sharper question. The configuration looks strong on development data. Does it survive contact with queries it has never seen?

\section{Protected Confirmation and Full-Scale Diagnosis}

\begin{figure*}[!t]
\centering
\includegraphics[width=0.96\textwidth]{figures/rebuilt/fig2_a5_confirmation.pdf}
\caption{\textbf{Final-872 comparison of two complete frozen configurations.} (a)~Absolute cross-domain Recall@100 and nDCG@100. (b)~Paired selected-minus-comparator differences with 95\% percentile intervals from 10,000 paired bootstrap resamples. (c)~Recall@100 query-level wins, ties, and losses. The comparison supports the complete selected system, not a representation-only causal effect.}
\label{fig:confirmation}
\end{figure*}

\subsection{Protected Held-Out Confirmation}

\begin{table}[!t]
\caption{Complete frozen systems.}
\label{tab:systems}
\centering
\footnotesize
\begin{tabularx}{\columnwidth}{@{}lX@{}}
\toprule
System & Frozen specification \\
\midrule
Research & \selectedSpec \\
Comparator & \comparatorSpec \\
\bottomrule
\end{tabularx}
\end{table}

It does. Four frozen profiles enter one Selection-125 evaluation with complete rankings; cross-domain selection statistics use 90 queries with positive cross-domain relations. The preregistered rule selects the research system; the research system and static comparator are fixed at that close, with complete specifications in Table~\ref{tab:systems}. After Selection-125, both complete systems---model, representation, prompt, program, runtime---are locked before Final-872.

Both systems run all 872 queries without failure; deterministic replay is verified. The selected configuration lifts Recall@100 from 0.331 to 0.442, a paired difference of 0.111 (95\% CI [0.102, 0.120]); it wins 619 queries, ties 158, and loses 95 (Figure~\ref{fig:confirmation}). nDCG@100 rises from 0.279 to 0.366 (0.086, [0.079, 0.094]), and nDCG@10 from 0.234 to 0.297 (0.064, \ndcgtenCI).

Development and selection were closed before these outcomes were observed, allowing the comparison to carry weight. It is the paper's strongest confirmatory evidence and is deliberately system-level: the design fixes representation with every other component, so the win belongs to the complete configuration, not representation alone.

\subsection{Full-Benchmark Candidate Exposure}

The selected system wins its comparison decisively. On the strict cross-domain slice, it still misses most of what it should find---a confirmed system is the right place to ask why. We keep the winner fixed and run it unchanged over the full benchmark: 45,336 target families, 1,247 queries, Top-200 per query. The run materializes 188,944 chunks and 249,400 ranked family rows with full coverage, no failures, and no duplicate-family errors. This strict cross-domain population differs from Final-872, so we report it separately from the confirmation.

\begin{figure*}[!t]
\centering
\includegraphics[width=0.96\textwidth]{figures/rebuilt/fig3_a7_diagnosis.pdf}
\caption{\textbf{Candidate exposure and within-pool ordering bound.} Post-confirmatory diagnosis of the fixed full-benchmark Top-200 pool for strict cross-domain relations. (a)~Relevant-family incidence exposure: 796 by rank 100, 332 first at ranks 101--200, and 4,065 absent. (b)~The analytical perfect-ordering bound within the immutable pool raises macro Recall@100 from 0.188 to 0.260, a headroom of 0.072; it is not reranker performance. Incidences and macro Recall use different aggregation units.}
\label{fig:diagnosis}
\end{figure*}

The pool licenses one bounded question: how much remaining error could a perfect reordering of existing candidates recover? Under the strict cross-domain definition, 905 judged queries hold 5,193 relevant-family incidences. Observed macro Recall@100 is 0.188; at depth 200 it is 0.260.

Figure~\ref{fig:diagnosis} separates exposure from ordering, and the split is stark. Of 5,193 relevant incidences, 796 surface by rank 100 and 332 more by rank 200---and 4,065 never appear at all. Seventy-eight percent of the evidence is not late; it is absent. The query view agrees: only 67 of 905 queries have every relevant family inside the top 100, while 455 expose none by rank 200. Reorder the pool perfectly and macro Recall@100 climbs from 0.188 to 0.260---a headroom of 0.072, and the entire prize available to any reranker. Everything above it was never retrieved. For a first-stage configuration this strong, the binding constraint is exposure, not ordering.

\section{Discussion and Conclusion}

Three results, one lesson: the unit that deserves a name in cross-domain patent retrieval is the retriever--representation configuration, not the retriever alone. The model scores whatever the construction, view selection, segmentation, and family aggregation make visible; these choices deserve the precision of a checkpoint name or retrieval depth. But precision is not leverage. At the coarse, field level these choices do not reorder retrievers---the reordering that looks real dissolves under resampling, while the retrievers stay firmly apart. Representation specifies the evidence; the retriever decides what to do with it.

The protected comparison moves the burden of proof from development behavior to a complete system that holds up on held-out queries with development closed. We keep the two levels apart on purpose: the transfer matrix shapes how we read the result, but only the held-out comparison is allowed to confirm.

The confirmed system points to where the next gain lives---and it is not the ranker. Four in five relevant families never reach the candidate pool, so reranking is capped at a headroom of 0.072 while candidate expansion still has most of the recall to win. The two interventions act on different failures: expansion changes what is exposed, reranking changes how exposed evidence is ordered. We do not rank them against each other here; we say which one the numbers make urgent.

Everything about representation and transfer is development evidence from a single benchmark, and no individual transfer cell carries confirmatory weight. The full-scale diagnosis examines one selected dense configuration. PatEmbed-large carries a research/non-commercial license. And citation-based patent retrieval is an information-retrieval task, not a legal reading of novelty, validity, infringement, or freedom to operate; nothing here should be read as the latter.

A retrieval score reflects both the model that scores evidence and the construction that decides what evidence exists. Once the first-stage configuration is fixed, a further constraint appears: which relevant families enter the pool. Together they set the reporting unit and the honest limit of downstream ranking.

\bibliographystyle{IEEEtran}
\bibliography{references/references}

\end{document}
